%! AP Statistics — Unit 9, Lesson 9.1 — Guided Notes (Answer Key)
\documentclass[10pt]{article}
\usepackage{apstats-article}
\usepackage{apstats-key}

\begin{document}

\pageheader{Unit 9, Lesson 9.1}{Guided Notes --- Answer Key}
\namedateperiod

\begin{objectivebox}
By the end of this lesson, I will be able to:
\begin{itemize}
    \item[$\square$] \ansline{Explain why the sample slope $b$ is a statistic that varies from sample to sample.}
    \item[$\square$] \ansline{State the null hypothesis for slope and explain what $\beta = 0$ means in context.}
    \item[$\square$] \ansline{Identify the central question of Unit 9: is the observed slope statistically significant evidence that $\beta \ne 0$?}
\end{itemize}
\end{objectivebox}

\begin{vocabbox}
\termblanklong{Population slope ($\beta$)}
\ansline{The true rate of change in the linear relationship between two quantitative variables in the population; a fixed (unknown) parameter.}

\termblanklong{Sample slope ($b$)}
\ansline{The slope of the LSRL computed from sample data; a statistic that varies from sample to sample.}

\termblanklong{Sampling variability}
\ansline{The natural variation in a statistic (such as $b$) across different random samples drawn from the same population.}

\termblanklong{Null hypothesis for slope}
\ansline{$H_0\colon \beta = 0$ --- the population has no linear relationship between $x$ and $y$; any observed non-zero $b$ is due to chance.}
\end{vocabbox}

\begin{hookbox}
\textbf{Discussion:} Two scatter plots are shown on the board. Both have a
positive LSRL with slope $b > 0$. One has tightly clustered points; the other
has widely scattered points. Which trend are you more confident is real? Why?

\begin{teachernote}
    The tightly clustered plot provides stronger evidence of a real relationship.
    Large scatter means $b$ could vary widely from sample to sample, so a non-zero
    $b$ from scattered data is less surprising under $H_0\colon \beta = 0$.
    This discussion primes students for the role of variability in inference.
\end{teachernote}
\end{hookbox}

\begin{notesbox}{The Sample Slope Is a Statistic}
In Units 2--3 we fit an LSRL $\hat{y} = a + bx$ to sample data and computed
the slope $b$. The slope $b$ is a \ans{statistic}, not a parameter.

\vspace{0.1in}
\textbf{Key idea:} Different random samples from the same population produce
\ans{different (varying)} values of $b$.

\vspace{0.1in}
The true (unknown) rate of change in the population is called the
\textbf{population slope}, denoted \ans{$\beta$} (``beta'').

\vspace{0.1in}
\renewcommand{\arraystretch}{1.6}
\begin{tabularx}{\linewidth}{Xcc}
    & \textbf{Symbol} & \textbf{Parameter or Statistic?} \\
    \hline
    True population slope       & $\beta$ & \ans{Parameter} \\
    Estimated slope from sample & $b$     & \ans{Statistic} \\
\end{tabularx}
\end{notesbox}

\begin{notesbox}{The Central Question of Unit 9}
When we see a non-zero slope $b$ in sample data, two explanations are possible:

\begin{enumerate}
  \item The true population slope $\beta$ \ans{is not zero ($\beta \ne 0$)} --- the trend is real.
  \item The true population slope $\beta = $ \ans{0} --- there is no real
        linear relationship, and the non-zero $b$ is due to \ans{random chance / sampling variability}.
\end{enumerate}

\vspace{0.1in}
\textbf{The null hypothesis for slope:} $H_0\colon \beta = $ \ans{0}

\vspace{0.08in}
This says: in the population, there is \ans{no linear relationship}
between $x$ and $y$. Any non-zero $b$ we observe is just \ans{random variation}.

\vspace{0.1in}
\textbf{Unit 9 asks:} Is the sample slope $b$ we observed \ans{statistically significant}
evidence that $\beta \ne 0$ in the population?
\end{notesbox}

\begin{notesbox}{Why Scatter Matters}
Two factors influence how surprising a non-zero $b$ is under $H_0\colon \beta = 0$:

\begin{enumerate}
  \item \textbf{Sample size ($n$):} A slope of $b = 3.4$ from $n = 200$ observations
        is \ans{stronger} evidence than the same slope from $n = 10$ observations.
  \item \textbf{Scatter around the line:} The more tightly points cluster around the
        LSRL, the \ans{more reliable / precise} our estimate of $\beta$.
\end{enumerate}

\vspace{0.1in}
This is why we can't just look at $b$ and declare the relationship real --- we need
\ans{a formal significance test (inference)} to quantify how likely $b$ is under $H_0$.
\end{notesbox}

\begin{notesbox}{Unit 9 Arc Overview}
\renewcommand{\arraystretch}{1.5}
\begin{tabularx}{\linewidth}{lX}
    \textbf{Lesson 9.1} & Framing: what question does Unit 9 ask? \\
    \textbf{Lesson 9.2} & Sampling distribution of $b$; conditions for inference \\
    \textbf{Lesson 9.3} & Confidence interval for \ans{$\beta$ (slope)} \\
    \textbf{Lesson 9.4} & Significance test for \ans{$\beta$} ($t$-test for slope) \\
    \textbf{Lesson 9.5} & Reading \ans{computer} output for regression \\
    \textbf{Lesson 9.6} & Synthesis: \ans{choosing and applying inference procedures for slopes} \\
\end{tabularx}
\end{notesbox}

\begin{practicebox}
\textbf{Guided Practice:} A study collects data on daily high temperature ($x$, in $^\circ$F)
and the number of ice cream cones sold ($y$) at a beachside stand for 20 days. The LSRL is
$\hat{y} = -82.4 + 3.4x$ with $r = 0.87$.

\begin{enumerate}
  \item Interpret the slope in context.
        \ansline{For each additional degree (Fahrenheit) in daily high temperature, the predicted number of ice cream cones sold increases by 3.4.}

  \item Suppose the true population slope were $\beta = 0$. What would that mean
        about the relationship between temperature and ice cream sales?
        \ansline{There is no real linear relationship between temperature and ice cream sales in the population. Any non-zero sample slope $b$ would be due to random variation in the sample.}

  \item Does $r = 0.87$ alone prove that $\beta \ne 0$? Explain.
        \ansline{No. $r = 0.87$ is a sample statistic; it could be non-zero by chance even if $\beta = 0$. We need a significance test to determine how likely a sample correlation this large is under $H_0\colon \beta = 0$.}
\end{enumerate}
\end{practicebox}

\end{document}
